\documentclass[11pt]{article}
\usepackage[margin=1in]{geometry}
\usepackage{amsmath,amssymb,booktabs}
\usepackage[hidelinks]{hyperref}
\title{An explicit mediator origin for the attractive scalar phase}
\author{Zimmerman programme: CPB-UV-2026-10-16}
\date{16 October 2026}
\begin{document}
\maketitle
\section{What is established}
A canonical real mediator can generate the preceding host's attractive
quartic and stabilizing sextic terms in a renormalizable two-scalar
Lagrangian. An explicit one-parameter family has a unique globally stable
classical vacuum, an exactly characterized homogeneous charged coexistence
phase, and a controlled heavy-mediator limit at tree level. The matching
also changes the host's amplitude kinetic term. These results close the
specified classical mediator-potential question, not the full theory.
The response switching function, quantum thresholds, charge production,
formation, gravity and dark-energy normalization remain unresolved.

Complex charged scalars with real mediators are established soliton models,
including the Friedberg--Lee--Sirlin family \cite{uv:FLS}. Attractive cubic
couplings and quantum corrections have been investigated in related
models \cite{uv:attractive,uv:quantum}. These references establish related
context only; their stability results are not transferred to this potential.
No global novelty claim is made.

\section{Canonical particle model and exact vacuum bound}
Set the host vacuum mass $m_h=1$ as the mass unit and let $M\ge1$ be
the mediator mass in that unit. With $I=|\Psi|^2$ and real $\chi$, take
\begin{align}
 \mathcal L&=-|\partial\Psi|^2-\tfrac12(\partial\chi)^2-V(I,\chi),\\
 V&=I+L I^2+2M\chi I+U(\chi),\qquad L=2-\frac1M,\\
 U(\chi)&=\frac{M^2}{2}\chi^2-\frac M8\chi^3+\frac14\chi^4.
 \label{uv:potential}
\end{align}
All interactions have canonical mass dimension at most four; this is a
renormalizable classical particle model. This does not assert an
all-energy ultraviolet completion or quantum vacuum stability. Dimensionful
coefficients are restored by multiplying mass-dimension-one quantities by
$m_h$ and dimension-two quantities by $m_h^2$; in particular $L$ means
$2-m_h/M_{\rm physical}$.

For $\chi\ge0$ every $I$ term is nonnegative, while
\begin{equation}
 U(\chi)=\chi^2\left[\frac14(\chi-M/4)^2+\frac{31M^2}{64}\right]>0
 \quad(\chi\ne0).
\end{equation}
For $\chi=-x<0$, $U(-x)=M^2x^2/2+Mx^3/8+x^4/4>0$. If
$x\le1/(2M)$, the coefficient $1-2Mx$ of $I$ is nonnegative.
Otherwise minimize exactly over $I\ge0$:
\begin{align}
 I_{\rm min}&=(2Mx-1)/(2L),\\
 4L V_{\rm min}&=Mx\left[\frac L2x^2-2x+4\right]+Lx^4-1.
\end{align}
Since $1\le L<2$,
\begin{equation}
 \frac L2x^2-2x+4=\frac L2(x-2/L)^2+4-2/L\ge2.
\end{equation}
Thus $4LV_{\rm min}\ge2Mx+Lx^4-1>0$ for $x>1/(2M)$.
All field values have been covered: $V\ge0$ with equality only at
$I=\chi=0$, for every $M\ge1$. This is an algebraic global statement,
not an inference from numerical sampling.

Moreover
\begin{equation}
 U''(\chi)=M^2-\tfrac34M\chi+3\chi^2
 =3(\chi-M/8)^2+\frac{61M^2}{64}>0.
\end{equation}
The stationary mediator equation
$U'(\chi)+2MI=0$ therefore has one real solution for every $I\ge0$.
There is no branch ambiguity in its tree-level local elimination.

\section{Tree matching includes gradients}
For general coefficients
$V=m_h^2 I+\lambda_\Psi I^2+a\chi I+M^2\chi^2/2+
\mu\chi^3/3+\lambda_\chi\chi^4/4$, the local stationary branch is
\begin{equation}
 \chi_*(I)=-\frac a{M^2}I-\frac{\mu a^2}{M^6}I^2+O(I^3).
\end{equation}
Substitution yields
\begin{equation}
 V_{\rm eff}=m_h^2 I+
 \left(\lambda_\Psi-\frac{a^2}{2M^2}\right)I^2
 -\frac{\mu a^3}{3M^6}I^3+O(I^4).
\end{equation}
For \eqref{uv:potential}, $a=2M$, $\mu=-3M/8$ and
$\lambda_\chi=1$, hence
\begin{equation}
 V_{\rm eff}=I-\frac{I^2}{M}+\frac{I^3}{M^2}+O(I^4),
 \qquad \chi_*=-2I/M+\frac{3I^2}{2M^3}+O(I^3).
 \label{uv:sextic}
\end{equation}
At $M=1$ the first three terms are the previous illustrative sextic
potential; higher terms need not be negligible there.

Eliminating a propagating mediator also produces derivative operators:
\begin{equation}
 -\tfrac12(\partial\chi_*)^2
 =-\frac{2}{M^2}(\partial I)^2+\cdots.
\end{equation}
For $\Psi=f e^{i\theta}$ the exact adiabatic amplitude kinetic factor is
\begin{equation}
 Z_f=1+2I\left(\frac{2M}{U''(\chi_*)}\right)^2
 =1+\frac{8I}{M^2}+\cdots.
 \label{uv:kinetic}
\end{equation}
The temporal phase term is $-I(\partial\theta)^2$. These expressions
assume sufficiently slow backgrounds; integrating out $\chi$ dynamically
also generates higher derivatives. Therefore accurate potential matching
alone does not justify importing the old wall, radius or curvature results.

\section{Exact charged coexistence and its finite check}
For a homogeneous charged phase, $\Psi=\sqrt I e^{i\omega t}$,
the quantity determining coexistence is
$\omega_*^2=\min_{I>0,\chi}V/I$. At $\chi=-x$, minimizing over $I$
gives $I_*=\sqrt{U(-x)/L}$ and
\begin{equation}
 R(x)=1-2Mx+2\sqrt{L U(-x)}.
\end{equation}
The derivative is negative near $x=0$ and positive at large $x$.
Its unique positive stationary point is the positive root of
\begin{multline}
 Lx^4+\frac{3LM}{4}x^3+
 \left(\frac{137L}{64}-1\right)M^2x^2\\
 +\left(\frac{3L}{4}-\frac12\right)M^3x-M^3=0.
 \label{uv:root}
\end{multline}
All nonconstant coefficients are positive for $L\ge1$; the polynomial
is strictly increasing for $x>0$, from a negative value to infinity.
It follows that this root is the global minimum of $R$ on the negative
mediator branch. Positive $\chi$ cannot improve on $V/I\ge1$.
The vacuum bound and the negative initial slope give
$0<\omega_*^2<1$. The density is $n_*=2\omega_* I_*$ and the pressure
is zero: $V=\omega_*^2 I_*$. These are exact homogeneous classical
results, not a new finite-size two-field droplet solution.

The sextic truncation predicts $I_0=M/2$ and $\omega_0^2=3/4$.
Solving \eqref{uv:root} and the exact mediator branch gives
\begin{center}
\begin{tabular}{rrrrr}
\toprule
$M/m_h$ & $I_*$ & $\omega_*^2$ & potential error & $Z_f$\\
\midrule
1 & 0.530967 & 0.777373 & $-3.398\%$ & $1.574$\\
10 & 4.351594 & 0.774215 & $-2.585\%$ & $1.296$\\
100 & 48.989310 & 0.753423 & $-0.441\%$ & $1.039$\\
1000 & 498.929051 & 0.750358 & $-0.0475\%$ & $1.004$\\
\bottomrule
\end{tabular}
\end{center}
The error is $(V_{\rm sextic}-V_{\rm exact})/V_{\rm exact}$ evaluated
at the exact phase density. The first two fail the chosen one-percent
potential comparison, retained as false flags. The $M=10$ density shift
is approximately $-13\%$ despite its smaller potential error; derivatives
and stationary states require their own check. Exact amplitude/mediator
mass matrices at fixed $\omega_*$ have positive eigenvalues at the four
tested points; their lower eigenvalues are approximately
$0.3392,0.7355,0.9620,0.9960$. This is not a finite-body or full dynamical
stability certificate.

In the large-$M$ limit, \eqref{uv:root} divided by $M^3$ tends $x-1=0$.
Thus $x_*\to1$, $I_*/M\to1/2$, $\chi_*/M\to0$, $\omega_*^2\to3/4$
and $Z_f\to1$. With $\Psi=\sqrt M\,\psi$, the leading effective action
is $M$ times the previous unit-coefficient sextic action. This explains
the classical shape limit; the conserved charge and energy scale with $M$,
so the old fixed inventories must not be silently retained.

\section{What still prevents a full particle bridge}
The attractive quartic is the difference
$\lambda_\Psi-a^2/(2M^2)=-1/M$ of two order-one terms.
Its selected fractional cancellation is $1/(2M)$: $0.5\%$ at $M=100$
and $0.05\%$ at $M=1000$. Increasing the hierarchy improves the tree
expansion but requires increasingly accurate matching of that difference.
No symmetry or renormalization-group argument protecting it has been
provided. Loop thresholds and the independent mass, tadpole and vacuum
counterterms must be calculated before claiming quantum closure. No
scheme-dependent partial loop term is used as a vacuum prediction here.

The mediator potential also does not derive the previous phase-response
function $W(f,p)$ or the nodal medium. A new two-field finite-body solution
with its gradients and full response is needed before transferring any
stationarity or radial-curvature result. Charge creation/separation,
thermal populations, abundance, an energy receiver and source-dependent
formation remain open. The exterior-force calibration cost persists.
This real mediator is not automatically the original source $X$, its
charge-even bridge $\sigma$, or the auxiliary vacuum variable $Q$.
The vacuum magnitude, flux/history residual, acceleration normalization,
relativistic gravity and cosmological/empirical obligations remain open.

\section{Evidence and review}
\texttt{mediator\_phase\_matching\_run\_002} contains 37 exact/finite
controls: vacuum algebra, elimination and root identities, four exact
coexistence comparisons, sixteen stationary-branch samples and four
homogeneous amplitude-matrix tests. Run 001 completed the same core
calculation; run 002 adds the kinetic-factor output. Potential-validity
flags are scientific results, not forced passing assertions. The current
manifest passed input/output freshness validation. This is a self-review
with the explicit proof above, not an independent-referee or quantum
stability certificate.

\begin{thebibliography}{9}
\bibitem{uv:FLS} J. Heeck and M. Sokhashvili,
\emph{Revisiting the Friedberg-Lee-Sirlin soliton model},
arXiv:2303.09566v2, \url{https://arxiv.org/abs/2303.09566v2}.
\bibitem{uv:attractive} Y. Hamada, K. Kawana, T. Kim and P. Lu,
\emph{Q-Balls in the presence of attractive force},
arXiv:2407.11115v2, \url{https://arxiv.org/abs/2407.11115v2}.
\bibitem{uv:quantum} Y.-X. Su, Q.-X. Xie and S.-Y. Zhou,
\emph{Quantum-Corrected Q-balls in the Friedberg-Lee-Sirlin Model},
arXiv:2605.25243, \url{https://arxiv.org/abs/2605.25243}.
Only official primary abstracts/metadata of these three contextual sources
were checked here. Their detailed theorems and numerical claims are not
used as proof of this candidate.
\end{thebibliography}
\end{document}
